\documentclass[11pt,a4paper]{article}
\usepackage[T1]{fontenc}\usepackage[utf8]{inputenc}\usepackage[provide=*,italian]{babel}
\usepackage{lmodern,microtype,geometry}\geometry{margin=2.5cm}
\usepackage{amsmath,amssymb,amsthm,mathtools,booktabs,array,longtable,xcolor,enumitem,tikz}
\usetikzlibrary{arrows.meta,positioning,calc}
\usepackage[most]{tcolorbox}\usepackage{listings,hyperref,fancyhdr}
\definecolor{sapienzared}{RGB}{130,0,36}\definecolor{softred}{RGB}{249,241,243}
\definecolor{softblue}{RGB}{239,246,252}\definecolor{softgray}{RGB}{245,245,245}\definecolor{codegreen}{RGB}{25,110,65}
\hypersetup{colorlinks=true,linkcolor=sapienzared,urlcolor=sapienzared,pdftitle={Laboratorio Python e NetworkX},pdfauthor={Giampaolo Liuzzi}}
\pagestyle{fancy}\fancyhf{}\fancyhead[L]{\small Ottimizzazione su Reti}\fancyhead[R]{\small Capitolo 9}\fancyfoot[C]{\thepage}\renewcommand{\headrulewidth}{0.4pt}
\newtheorem{definition}{Definizione}[section]\newtheorem{proposition}[definition]{Proposizione}
\newtcolorbox{learningbox}{colback=softred,colframe=sapienzared,title=Obiettivi di apprendimento,fonttitle=\bfseries,boxrule=0.8pt,arc=1.5mm}
\newtcolorbox{examplebox}[1]{colback=softblue,colframe=blue!55!black,title=#1,fonttitle=\bfseries,boxrule=0.7pt,arc=1.5mm}
\newtcolorbox{notationbox}{colback=softgray,colframe=black!55,title=Convenzioni del progetto,fonttitle=\bfseries,boxrule=0.7pt,arc=1.5mm}
\newtcolorbox{warningbox}{colback=yellow!9,colframe=orange!65!black,title=Attenzione,fonttitle=\bfseries,boxrule=0.7pt,arc=1.5mm}
\newtcolorbox{networkxbox}{colback=softgray,colframe=black!55,title=Python/NetworkX,fonttitle=\bfseries,boxrule=0.7pt,arc=1.5mm}
\lstdefinestyle{python}{language=Python,basicstyle=\ttfamily\small,keywordstyle=\color{blue!65!black}\bfseries,commentstyle=\color{codegreen},stringstyle=\color{sapienzared},backgroundcolor=\color{softgray},frame=single,rulecolor=\color{black!25},showstringspaces=false,breaklines=true,columns=fullflexible}

\begin{document}\hypersetup{pageanchor=false}
\begin{titlepage}\centering\vspace*{1.6cm}
{\color{sapienzared}\rule{\textwidth}{1.2pt}}\\[1.1cm]{\Large Ottimizzazione su Reti}\\[0.35cm]
{\huge\bfseries Capitolo 9}\\[0.45cm]{\LARGE\bfseries Laboratorio Python e NetworkX}\\[1.1cm]
{\color{sapienzared}\rule{0.72\textwidth}{0.8pt}}\\[1.1cm]
{\large Corso di Laurea Magistrale in Ingegneria Gestionale}\\[0.25cm]{\large Sapienza Università di Roma}\\[1.2cm]
{\large Giampaolo Liuzzi}\\[0.25cm]{\normalsize Dipartimento di Ingegneria Informatica, Automatica e Gestionale}
\vfill{\small Versione preliminare per uso didattico}\end{titlepage}
\hypersetup{pageanchor=true}\pagenumbering{roman}\tableofcontents\newpage\pagenumbering{arabic}

\begin{learningbox}
Al termine del capitolo lo studente dovrebbe essere in grado di:
\begin{itemize}[nosep]
\item rappresentare in modo uniforme grafi e istanze di ottimizzazione;
\item separare dati, algoritmo, verifica e visualizzazione;
\item controllare bilanci, capacità, costi e certificati di ottimalità;
\item usare correttamente le principali funzioni NetworkX;
\item realizzare visualizzazioni leggibili di flussi, cammini, alberi e tagli;
\item scrivere test automatici per piccole istanze;
\item organizzare notebook e codice in una repository riproducibile.
\end{itemize}
\end{learningbox}

\section{Obiettivi del laboratorio}
Il laboratorio computazionale non serve soltanto a ottenere un numero. Una buona soluzione deve rendere espliciti:
\begin{enumerate}
\item il modello matematico rappresentato dal grafo;
\item le convenzioni usate per nodi, archi e attributi;
\item l'algoritmo o la funzione di libreria impiegata;
\item i controlli indipendenti sul risultato;
\item un certificato di ottimalità quando disponibile.
\end{enumerate}

NetworkX offre strutture dati e algoritmi affidabili, ma non verifica se il grafo costruito descrive correttamente il problema reale. Modellazione e validazione restano responsabilità dell'utente.

\section{Ambiente di lavoro}
Per gli esempi sono sufficienti Python, NetworkX e Matplotlib; NumPy e pandas sono utili per matrici e tabelle.
\begin{lstlisting}[style=python,caption={Importazioni comuni}]
import networkx as nx
import matplotlib.pyplot as plt
import numpy as np
import pandas as pd
\end{lstlisting}

È opportuno registrare le versioni nel file delle dipendenze e usare un ambiente virtuale. I notebook devono poter essere eseguiti dall'inizio alla fine senza operazioni manuali nascoste.

\begin{networkxbox}
\texttt{nx.Graph} rappresenta grafi non orientati, \texttt{nx.DiGraph} grafi orientati. \texttt{MultiGraph} e \texttt{MultiDiGraph} ammettono archi paralleli, ma non tutte le funzioni di ottimizzazione li accettano.
\end{networkxbox}

\section{Convenzioni comuni per i dati}
Adottiamo nomi di attributi coerenti con NetworkX:
\begin{center}\begin{tabular}{lll}\toprule
Oggetto&Attributo&Significato\\\midrule
nodo&\texttt{demand}&$b_i$: positivo domanda, negativo offerta\\
arco&\texttt{weight}&costo, lunghezza o durata\\
arco&\texttt{capacity}&capacità superiore\\
arco&\texttt{lower}&limite inferiore, gestito esplicitamente\\
arco&\texttt{flow}&flusso candidato o calcolato\\
\bottomrule\end{tabular}\end{center}

\begin{notationbox}
La matrice di incidenza ha $-1$ all'origine e $+1$ alla destinazione. Pertanto
\[
Ax=b,\qquad b_i>0\text{ domanda},\qquad b_i<0\text{ offerta}.
\]
L'attributo \texttt{demand} coincide direttamente con $b_i$: non deve essere cambiato di segno.
\end{notationbox}

\subsection{Liste di archi e tabelle}
Una piccola istanza può essere definita con tuple; per dati esterni è preferibile una tabella CSV con colonne stabili.
\begin{lstlisting}[style=python,caption={Costruzione da una tabella di archi}]
edges = pd.DataFrame([
    {"tail": 1, "head": 2, "cost": 4, "capacity": 10},
    {"tail": 1, "head": 3, "cost": 2, "capacity": 7},
    {"tail": 2, "head": 3, "cost": 1, "capacity": 5},
])

G = nx.DiGraph()
for row in edges.itertuples(index=False):
    G.add_edge(row.tail, row.head,
               weight=row.cost, capacity=row.capacity)
\end{lstlisting}

\section{Validazione dell'istanza}
I controlli devono precedere la chiamata all'algoritmo. Verifichiamo almeno:
\begin{itemize}
\item tipo di grafo e presenza di nodi e archi;
\item attributi obbligatori e valori numerici finiti;
\item capacità non negative e limiti inferiori non superiori alle capacità;
\item somma delle domande uguale a zero per il flusso a costo minimo;
\item esistenza dei nodi sorgente e destinazione;
\item aciclicità per algoritmi su DAG e CPM;
\item non negatività dei pesi prima di usare Dijkstra.
\end{itemize}

\begin{lstlisting}[style=python,caption={Validatore essenziale per una rete di flusso}]
def validate_flow_instance(G):
    if not G.is_directed():
        raise ValueError("La rete deve essere orientata")

    total_demand = sum(G.nodes[v].get("demand", 0) for v in G)
    if total_demand != 0:
        raise ValueError("La somma delle domande deve essere zero")

    for i, j, data in G.edges(data=True):
        cap = data.get("capacity", float("inf"))
        low = data.get("lower", 0)
        if low < 0 or cap < low:
            raise ValueError(f"Limiti non validi su {(i, j)}")
\end{lstlisting}

\section{Verifica indipendente delle soluzioni}

\subsection{Bilanci e capacità}
Un flusso candidato va controllato senza affidarsi al fatto che la funzione sia terminata correttamente.
\begin{lstlisting}[style=python,caption={Verifica di un flusso nella convenzione $Ax=b$}]
def check_flow(G, flow, tol=1e-9):
    balance = {v: 0.0 for v in G}
    cost = 0.0

    for i, j, data in G.edges(data=True):
        x = flow[i][j]
        low = data.get("lower", 0.0)
        cap = data.get("capacity", float("inf"))
        if x < low - tol or x > cap + tol:
            return False, None
        balance[i] -= x
        balance[j] += x
        cost += data.get("weight", 0.0) * x

    ok = all(abs(balance[v] - G.nodes[v].get("demand", 0.0))
             <= tol for v in G)
    return ok, cost
\end{lstlisting}

\subsection{Certificati specifici}
\begin{center}\begin{tabular}{ll}\toprule
Problema&Controllo o certificato\\\midrule
flusso a costo minimo&bilanci, limiti, costo e costi ridotti\\
massimo flusso&valore del flusso uguale alla capacità del taglio\\
cammino minimo&costo del cammino e disuguaglianze di Bellman\\
albero&connessione, aciclicità e $n-1$ spigoli\\
CPM&ricorrenze avanti/indietro e slittamenti non negativi\\
\bottomrule\end{tabular}\end{center}

\section{Mappa delle principali funzioni NetworkX}
\begin{center}\small
\begin{tabular}{lll}\toprule
Problema&Funzione&Attributo\\\midrule
componenti connesse&\texttt{connected\_components}&--\\
ordinamento topologico&\texttt{topological\_sort}&--\\
albero ricoprente minimo&\texttt{minimum\_spanning\_tree}&\texttt{weight}\\
flusso a costo minimo&\texttt{network\_simplex}&\texttt{demand}, \texttt{weight}, \texttt{capacity}\\
massimo flusso&\texttt{maximum\_flow}&\texttt{capacity}\\
taglio minimo&\texttt{minimum\_cut}&\texttt{capacity}\\
cammino minimo&\texttt{shortest\_path}&\texttt{weight}\\
Dijkstra&\texttt{single\_source\_dijkstra}&\texttt{weight}\\
Bellman-Ford&\texttt{single\_source\_bellman\_ford}&\texttt{weight}\\
cammino massimo in DAG&\texttt{dag\_longest\_path}&\texttt{weight}\\
\bottomrule\end{tabular}\end{center}

\begin{warningbox}
Il nome predefinito dell'attributo è spesso parte dell'interfaccia. Se i dati usano \texttt{cost} ma la funzione legge \texttt{weight}, bisogna passare il parametro appropriato o rinominare l'attributo. Un attributo ignorato può produrre un risultato formalmente valido ma riferito a un problema diverso.
\end{warningbox}

\section{Separare algoritmo e libreria}
Per ciascun argomento è utile affiancare:
\begin{enumerate}
\item una versione didattica, breve e tracciabile;
\item la funzione NetworkX per il confronto;
\item una funzione indipendente di verifica.
\end{enumerate}

La versione didattica deve restituire non solo la soluzione finale, ma anche informazioni sulle iterazioni: nodi etichettati, archi rilassati, cammini aumentanti, alberi di base o costi ridotti. La visualizzazione non deve essere incorporata nel nucleo algoritmico: è meglio restituire dati strutturati e disegnarli in una funzione separata.

\begin{examplebox}{Schema di una funzione didattica}
Una funzione per Dijkstra può restituire \texttt{dist}, \texttt{pred} e una lista \texttt{trace}. Ogni elemento di \texttt{trace} contiene il nodo reso permanente e una copia delle etichette dopo i rilassamenti. Lo stesso schema permette di costruire tabelle o animazioni senza modificare l'algoritmo.
\end{examplebox}

\section{Visualizzazione delle reti}
Una figura efficace deve distinguere struttura e soluzione:
\begin{itemize}
\item posizione dei nodi salvata e riutilizzata;
\item etichette degli archi coerenti, per esempio costo/capacità o flusso/capacità;
\item cammino, albero o taglio evidenziato con colore e spessore;
\item archi non rilevanti attenuati, senza eliminare informazione necessaria.
\end{itemize}

\begin{lstlisting}[style=python,caption={Evidenziare un cammino}]
def draw_path(G, path, weight="weight"):
    pos = nx.spring_layout(G, seed=7)
    path_edges = set(zip(path, path[1:]))
    colors = ["#820024" if e in path_edges else "#bbbbbb"
              for e in G.edges]
    widths = [2.8 if e in path_edges else 1.0 for e in G.edges]

    nx.draw(G, pos, with_labels=True, node_color="#f9f1f3",
            edge_color=colors, width=widths, arrows=True)
    labels = nx.get_edge_attributes(G, weight)
    nx.draw_networkx_edge_labels(G, pos, edge_labels=labels)
    plt.tight_layout()
\end{lstlisting}

Per un taglio minimo si colorano diversamente i nodi di $W$ e $\overline W$; per un flusso si usa l'etichetta $x_{ij}/u_{ij}$; per il CPM si evidenziano gli archi a slittamento nullo.

\section{Test automatici}
I test su istanze piccole proteggono da errori di segno e regressioni. Sono particolarmente utili proprietà che devono valere per ogni soluzione.
\begin{lstlisting}[style=python,caption={Esempi di test con pytest}]
def test_maxflow_equals_mincut(G):
    value, _ = nx.maximum_flow(G, "s", "t")
    cut_value, _ = nx.minimum_cut(G, "s", "t")
    assert value == cut_value

def test_shortest_path_cost(G):
    path = nx.shortest_path(G, "s", "t", weight="weight")
    cost = nx.path_weight(G, path, weight="weight")
    assert cost == nx.shortest_path_length(
        G, "s", "t", weight="weight"
    )

def test_tree_size(T):
    assert nx.is_tree(T)
    assert T.number_of_edges() == T.number_of_nodes() - 1
\end{lstlisting}

Con dati in virgola mobile si usa una tolleranza, per esempio \texttt{pytest.approx} o \texttt{math.isclose}. È utile includere anche casi non validi: domanda sbilanciata, capacità negativa, DAG con ciclo e Dijkstra con peso negativo.

\section{Struttura della repository GitHub}
Una possibile organizzazione è:
\begin{verbatim}
ottimizzazione-su-reti/
|-- README.md
|-- pyproject.toml
|-- src/network_optimization/
|   |-- validation.py
|   |-- visualization.py
|   |-- flows.py
|   |-- shortest_paths.py
|   `-- cpm.py
|-- notebooks/
|   |-- 01_grafi.ipynb
|   |-- 02_flusso_costo_minimo.ipynb
|   |-- 03_massimo_flusso.ipynb
|   |-- 04_cammini_minimi.ipynb
|   `-- 05_cpm.ipynb
|-- data/
|-- tests/
`-- exercises/
\end{verbatim}

Il \texttt{README} dovrebbe indicare installazione, convenzioni, ordine dei notebook e licenza. I notebook illustrano gli esempi; le funzioni riutilizzabili appartengono a \texttt{src}; i test non vanno nascosti nei notebook.

\section{Riproducibilità e stile}
\begin{itemize}
\item fissare il seme dei layout casuali;
\item non dipendere dall'ordine accidentale di nodi o archi;
\item evitare variabili globali e celle eseguite fuori ordine;
\item aggiungere docstring, type hints e nomi coerenti;
\item salvare dati di esempio in formati testuali semplici;
\item confrontare implementazione didattica e libreria senza assumere che restituiscano lo stesso cammino in caso di ottimi multipli.
\end{itemize}

\section{Errori frequenti}
\begin{enumerate}
\item Cambiare segno alle domande prima di passarle a NetworkX.
\item Usare \texttt{Graph} quando il modello richiede archi orientati.
\item Confondere \texttt{weight} e \texttt{capacity}.
\item Fidarsi del risultato senza verificare bilanci e vincoli.
\item Confrontare cammini o alberi invece dei valori quando esistono ottimi multipli.
\item Incorporare disegno, input e algoritmo nella stessa funzione.
\item Usare Dijkstra senza controllare i pesi.
\item Costruire il CPM su un grafo ciclico.
\item Non fissare il layout, ottenendo figure diverse a ogni esecuzione.
\end{enumerate}

\section{Esercizi di laboratorio}
\textbf{Esercizio 9.1 -- Importazione.} Definire un formato CSV per nodi e archi, importare una rete e validarne gli attributi.

\medskip\textbf{Esercizio 9.2 -- Flusso.} Risolvere un flusso a costo minimo con \texttt{network\_simplex} e verificarne indipendentemente bilanci e costo.

\medskip\textbf{Esercizio 9.3 -- Certificato.} Calcolare massimo flusso e taglio minimo, visualizzando la partizione e gli archi del taglio.

\medskip\textbf{Esercizio 9.4 -- Confronto.} Implementare Dijkstra didattico con traccia e confrontare valori e predecessori con NetworkX.

\medskip\textbf{Esercizio 9.5 -- CPM.} Produrre una tabella pandas con $ES$, $EF$, $LS$, $LF$ e slittamento di ogni attività.

\medskip\textbf{Esercizio 9.6 -- Test.} Scrivere almeno un test positivo e uno negativo per ciascun algoritmo.

\section{Sintesi del capitolo}
\begin{learningbox}\begin{itemize}[nosep]
\item Una convenzione comune rende interoperabili modelli, esempi e notebook.
\item NetworkX risolve il grafo fornito, non verifica la correttezza del modello.
\item Ogni soluzione dovrebbe essere accompagnata da controlli e certificati.
\item Algoritmo, validazione e visualizzazione vanno mantenuti separati.
\item I test automatici rendono affidabili le implementazioni didattiche.
\item La repository deve distinguere codice riutilizzabile, notebook, dati ed esercizi.
\end{itemize}\end{learningbox}

\paragraph{Passo successivo.} La struttura proposta costituisce la base naturale per realizzare la pagina GitHub del corso e pubblicare progressivamente notebook, istanze ed esercizi.
\end{document}
